\chapter*{บทคัดย่อ}

การวิจัยนี้มีวัตถุประสงค์เพื่อศึกษาผลกระทบของการเข้าร่วมโครงการเกษตรกรอัจฉริยะ (Smart Farmer) ที่มีต่อทักษะของเกษตรกร โดยใช้ข้อมูลจากเกษตรกรที่เข้าร่วมและไม่เข้าร่วมโครงการในช่วงปี พ.ศ. 2565–2566 การวิเคราะห์ใช้เทคนิคทางเศรษฐมิติร่วมกับวิธีการเรียนรู้ของเครื่อง (Machine Learning) ได้แก่ Propensity Score Matching (PSM) และ Difference-in-Differences (DID) เพื่อควบคุมอคติจากการคัดเลือกตัวอย่างและประเมินผลกระทบเชิงสาเหตุ

ผลการวิจัยพบว่า แม้เกษตรกรที่เข้าร่วมโครงการจะมีระดับทักษะสูงกว่ากลุ่มที่ไม่เข้าร่วมทั้งก่อนและหลังการดำเนินโครงการ แต่เมื่อควบคุมปัจจัยด้วยวิธี DID แล้ว ไม่พบว่าการเข้าร่วมโครงการส่งผลให้เกิดการพัฒนาทักษะที่เพิ่มขึ้นอย่างมีนัยสำคัญทางสถิติ นอกจากนี้ การวิเคราะห์ด้วย Machine Learning และ SHAP พบว่าปัจจัยสำคัญที่มีผลต่อการพัฒนาทักษะ ได้แก่ ระดับทักษะตั้งต้น ทักษะด้านเทคโนโลยี ระดับการศึกษา และโครงสร้างพื้นฐาน เช่น ระบบชลประทาน

ผลการศึกษาแสดงให้เห็นว่าผลลัพธ์ของโครงการมีความแตกต่างกันในแต่ละกลุ่มเกษตรกร (heterogeneity) และควรออกแบบหลักสูตรและนโยบายแบบจำเพาะกลุ่ม (targeted policy) เพื่อเพิ่มประสิทธิผลของโครงการในอนาคต

\vspace{0.5cm}
\textbf{คำสำคัญ:} เกษตรกรอัจฉริยะ, ทักษะเกษตรกร, การประเมินผลกระทบ, การจับคู่คะแนนความโน้มเอียง, การเรียนรู้ของเครื่อง